\section{Strictly aspherical endpoint functors}
\label{sec:C}

Throughout this section $Y$ is a small category with specified object
$y_0$, $P=P(Y,y_0)$ is the full path category of the question, and
$q:P\to Y$ is endpoint evaluation. No bound on path length is imposed.
For $z\in Y$ put $H_z=Y(q(-),z):P^{\op}\to\Set$ and let
$\pi_z:P/H_z\to P$ denote its element projection.

\begin{claim}{C.1}{Proved}{The integral obstruction for an outgoing arrow}
\label{C:outgoing}
Let $u:y_0\to z$. Define paths
\[
 d=(y_0\xrightarrow{u}z),\qquad
 p=(y_0\xrightarrow{1}y_0\xrightarrow{u}z).
\]
There are morphisms $f,g:p\to d$ in $P$ whose vertex maps are
$(0,1,1)$ and $(0,0,1)$ and whose vertical components are, respectively,
$(1,u,1_z)$ and $(1,1,1_z)$. Let
\[
 L=P(p,-)/(kf\sim kg):P\longrightarrow\Set,
\]
where at $x$ we impose this relation for every $k:d\to x$. Then
\[
 H_1\bigl(B\El_{P/H_z}(L\pi_z);\Z\bigr)
 \cong \Z[Y(z,z)]\big/
 \big\langle\gamma v\mid
 \gamma:y_0\to z,\ v:z\to y_0,\ vu=1_{y_0}\big\rangle.
 \tag{C.1.1}\label{eq:C-outgoing}
\]
The denominator is the subgroup spanned by the specified basis elements.
It equals the numerator only if $u$ is invertible.
\end{claim}
\begin{proof}
The first and second naturality squares for $f$ are
$u1=u1$ and $1_zu=1_zu$; those for $g$ are identities. Thus the
specified maps exist and both have endpoint component $1_z$.
Write $P_a=\Z[P(a,-)]\in\Ab^P$ and define
$\delta:P_d\to P_p$ by $k\mapsto kf-kg$. With
$K=\ker\delta$ and $J=\operatorname{im}\delta$ there are exact
sequences of covariant $\Z P$-modules
\[
 0\to K\to P_d\to J\to0,
 \qquad 0\to J\to P_p\to\Z[L]\to0.
\]
The two representables are projective. Tensoring with the right module
$\Z[H_z]$ and using co-Yoneda consequently identifies
\[
 \Tor^{\Z P}_1(\Z[H_z],\Z[L])
 =\ker\left(\Z[H_z(d)]/N\longrightarrow\Z[H_z(p)]\right),
\]
where $N$ is the image of
$\Z[H_z]\otimes_P K\to\Z[H_z(d)]$.
The displayed map is induced by $H_z(f)-H_z(g)=0$.
The two-sided bar resolution identifies this Tor group with the
homology in the statement: its basis consists of a composable chain
in $P$, an $L$-element at its initial object, and an $H_z$-element
at its terminal object. These are exactly the nerve chains of the
specified restricted action category. This verifies the variance and
the step of PRETOT A.1 used here.

Explicitly, $N$ is generated by
\[
 \gamma\cdot t=\sum_i n_i\gamma q(k_i),\qquad
 t=\sum_i n_i k_i\in K(x),\quad \gamma\in H_z(x).
\]
If $x$ has nonempty type of length $m$, any map $k:d\to x$
has vertex map $(0,m)$. The composite $kf$ has vertex map
$(0,m,m)$, whereas every $k'g$ has vertex map $(0,0,m)$.
These supports are disjoint. The equality $tf=tg$ therefore makes
each side zero separately. Within a fiber of $k\mapsto kf$ the
endpoint component $q(k)=q(kf)$ is constant. The coefficients in
that fiber sum to zero, so $\gamma\cdot t=0$.
This argument treats every $m>0$, including arbitrarily long paths
and all possible endomorphisms in $Y$.

The only empty path is $o$. A morphism $d\to o$ is exactly an
arrow $v:z\to y_0$ satisfying $vu=1_{y_0}$. Each such morphism
equalizes $f,g$: both composites have vertical components
$(1,1,v)$. Hence $K(o)=P_d(o)$ and its tensor image consists
exactly of the elements $\gamma v$ in \eqref{eq:C-outgoing}.
This proves the formula.

A subgroup spanned by a subset of the standard basis of a free
abelian group contains the basis vector $1_z$ only if that vector
belongs to the subset. Thus vanishing gives $\gamma v=1_z$ and
$vu=1_{y_0}$ for some $\gamma,v$. Then
$\gamma=\gamma vu=u$, so $u$ and $v$ are inverse.
\end{proof}

\begin{claim}{C.2}{Proved}{The obstruction for an incoming arrow}
\label{C:incoming}
For $v:z\to y_0$, put
\[
 d=(y_0\xleftarrow{v}z),\qquad
 p=(y_0\xleftarrow{v}z\xleftarrow{1_z}z).
\]
Use the same vertex maps for $f,g:p\to d$, with components
$(1,1_z,1_z)$ and $(1,v,1_z)$. For the corresponding quotient
$L=P(p,-)/(kf\sim kg)$ one has
\[
 H_1\bigl(B\El_{P/H_z}(L\pi_z);\Z\bigr)
 \cong \Z[Y(z,z)]/
 \langle\gamma v\mid\gamma:y_0\to z\rangle.
 \tag{C.2.1}
\]
Vanishing implies that $v$ has a left inverse.
\end{claim}
\begin{proof}
The naturality equations verify the displayed maps and again
$q(f)=q(g)=1_z$. The projective presentation and the disjoint-support
calculation of C.1 apply without change. There is exactly one arrow
$d\to o$, with last component $v$; it equalizes $f,g$.
This gives the denominator. Its containing $1_z$ gives
$\gamma v=1_z$, precisely a left inverse.
\end{proof}

\begin{claim}{C.3}{Proved}{The full walking-retraction counterexample}
\label{C:retraction}
For the walking retraction $R_0$, with objects $0,z$, arrows
$u:0\to z$, $r:z\to0$, relations $ru=1_0$, $e=ur\ne1_z$,
and basepoint $0$, both C.1 and C.2 give
\[
 H_1\bigl(B\El_{P/H_z}(L\pi_z);\Z\bigr)
   =\Z\{1_z,e\}/\Z e\cong\Z.
\]
Consequently $P(R_0,0)\to R_0$ is not strictly aspherical for
$W_\infty$, $W_{\Z}$, or any $W_n$ with $n\geq2$.
Moreover, as right $\Z P$-modules,
\[
 \Z[H_z]\cong\Z[P(-,o)]\oplus q^*S,
 \quad S(0)=0,\quad S(z)=\Z,
\]
where $e$ acts by zero on $S(z)$.
\end{claim}
\begin{proof}
The five arrows $1_0,1_z,u,r,e$ and the listed relations give the
entire category; composition is forced and can also be realized by
maps between the singleton $\{0\}$ and $\{0,1\}$, with $u$ the
inclusion and $r,e$ the constant maps. Here the sole incoming arrow
in the formulas is $r$ and the sole outgoing arrow is $u$, so the
denominator is $\Z e$.

The empty path is terminal since $0$ is terminal in $R_0$.
The natural basis element $u\circ(q(c)\to0)$ defines a split
copy of the constant module $\Z=\Z[P(-,o)]$ in $\Z[H_z]$.
At endpoint $z$ its complement is generated by $1_z-e$; at endpoint
$0$ it is zero. Precomposition by $e$ annihilates $1_z-e$.
This proves the claimed natural splitting.

Our independent program \texttt{CD\_finite\_checks.py} represents
arrows of $R_0$ as actual finite-set maps and builds paths directly
from the interval definition. For length at most two it verifies
25 objects and 585 arrows, closure under all composable pairs,
the two pairs $f,g$, the kernel-tensor image $\Z e$, and the
restricted action categories. They have 34 objects and respectively
803 and 819 arrows. An integer presentation obtained by deleting a
spanning forest from the nerve $1$-skeleton reduces by unimodular
unit pivots to $H_0=\Z^2$ and $H_1=\Z$ in each case; there is
no torsion. This is an exact integral computation, not only a
calculation modulo several primes.

We also read the supplied drafter's program in full. Its quotient
of $P(p,-)$ imposes exactly the pairs $kf\sim kg$; its element
objects carry both the covariant quotient label and the contravariant
$H_z$-label; its composition tests the requisite presheaf equation.
Its degree-two boundaries are the nerve boundaries, with identity
simplices retained consistently. Its integer-kernel calculation
computes the image of $\Z[H_z]\otimes K$ rather than merely the
rational kernel. Its stated finite computations therefore test
what Lead 2 claims. Neither program by itself proves a statement
about paths of all lengths; C.1--C.2 provide that proof.
\end{proof}

\begin{claim}{C.4}{Proved}{Classification of strict endpoint functors above degree one}
\label{C:strict-high}
For $W=W_\infty$, $W_{\Z}$, or any fixed $W_n$ with $n\geq2$,
the following are equivalent:
\begin{enumerate}
\item $q:P(Y,y_0)\to Y$ is strictly $W$-aspherical;
\item $q$ is totally $W$-aspherical;
\item $Y$ is a connected thin groupoid, equivalently $Y\simeq e$.
\end{enumerate}
\end{claim}
\begin{proof}
Strictness says that every $H_z$ is totally $W$-aspherical.
The restricted-action criteria verified in PRETOT A.1 and C.1
make the first homology of every restricted-action component zero
for each of the listed localizers. Thus C.1 makes every arrow
out of $y_0$ invertible. By C.2 an incoming arrow $v:z\to y_0$
has a left inverse $\gamma:y_0\to z$; the already proved outgoing
case makes $\gamma$ invertible, so $v=\gamma^{-1}$.

We check the multiplicity step from PRETOT E.7 also for the truncated
localizers. The empty path forms a full cosieve $\{o\}\subset P$.
Let $E=P/H_z$ and $E_o=\{o\}/H_z$, the discrete category
$Y(y_0,z)$. Take the $P$-action with two elements outside
$\{o\}$ and one element on it. Its restricted element category
is the double $E\amalg_{E_o}E$. At nerve level this is a pushout
along monomorphisms: any chain meeting a cosieve stays inside it.
The identity test makes $E$ nonempty connected and $H_1(BE)=0$.
If $E_o$ is nonempty, the double is connected and has $H_1=0$
by the same restricted-action criterion. Mayer--Vietoris then
forces the kernel of
$H_0(E_o)\to H_0(E)\oplus H_0(E)=\Z^2$ to vanish.
Thus $|Y(y_0,z)|\leq1$.

Every $E=P/H_z$ is nonempty. An object of it gives a path from
$y_0$ to some $x$ and an arrow $x\to z$, hence a zigzag from
$y_0$ to $z$. Thus $Y$ is connected. Proceeding along any such
zigzag, transport the next incident arrow across the isomorphism
already obtained from its preceding vertex to $y_0$. It becomes
an incoming or outgoing arrow at $y_0$ and is invertible. Every
object is therefore isomorphic to $y_0$, every arrow is invertible,
and the multiplicity bound makes every hom-set a singleton.

Conversely, in a connected thin groupoid there is exactly one
$\alpha_z:y_0\to z$. The object $(o,\alpha_z)$ is terminal
in $P/z$: collapse the interval and use the unique vertical arrows
to $y_0$. A category with terminal object has representable terminal
presheaf, whose left Kan extensions are representable. It is total
for every fundamental localizer. Thus (3) implies (2).
Finally a total functor is strict: for an arbitrary outgoing
$p:P\to U$ the canonical element comparison
$P/H_z\to U/p_!H_z$ is final, so totality of the source category
applies to its final image. This is the comprehensive-factorization
argument of PRETOT B.4, with exactly this variance. Hence (2)
implies (1).
\end{proof}

\begin{claim}{C.5}{Proved}{Low truncations of total presheaves}
\label{C:low-presheaves}
For an arbitrary small $U$ and presheaf $H:U^{\op}\to\Set$,
$H$ is totally $W_0$-aspherical exactly when $U/H$ is nonempty
and connected, and totally $W_1$-aspherical exactly when
$B(U/H)$ is nonempty and simply connected. The same criteria
hold for total categories.
\end{claim}
\begin{proof}
Necessity is the identity left-Kan-extension test. For any
$p:U\to U'$ the canonical functor
$j:U/H\to U'/p_!H$ is ordinary final, by the comprehensive
factorization: its outgoing commas are precisely the nonempty
connected components defining the pointwise left Kan extension.
An ordinary final functor induces a bijection on components and
surjects on each fundamental group. For the latter statement,
choose an object in each outgoing comma, replace every target edge
by its two chosen legs, and connect the alternative legs by a
finite zigzag inside the corresponding connected comma. The comma
triangles identify the resulting word in the target fundamental
groupoid with the original edge. Applying this to a loop gives a
lift modulo the connecting paths. Thus connectedness, and then
simple connectivity, transfer from source to target. This proves
sufficiency for every $p$. The category assertion is the case
$H=1_U$.
\end{proof}

\begin{claim}{C.6}{Known}{The ordinary slices are path spaces}
\label{C:slice-path}
For every $z\in Y$ there are specified weak equivalences
\[
 B(P/z)\simeq B\operatorname{Ch}(Y;y_0,z)
       \simeq\operatorname{Map}_{BY}(y_0,z).
 \tag{C.6.1}
\]
The second comparison is INDLOCAL C.7 with every arrow marked;
its proof is verified below at the step used.
\end{claim}
\begin{proof}
Let $j:\operatorname{Ch}(Y;y_0,z)\to P/z$ send a fixed-endpoint
path $c$ to $(c,1_z)$. Define
$r:P/z\to\operatorname{Ch}(Y;y_0,z)$ by appending the forward
edge $\beta:q(c)\to z$ to $(c,\beta)$.
On a morphism $k=(\phi,\eta):(c,\beta)\to(d,\beta')$, extend
$\phi$ by the last vertex and $\eta$ by $1_z$; the final square
is exactly $\beta'\eta_{\rm last}=\beta$.
The functor $rj$ has a natural transformation to the identity,
collapsing its added identity edge.

For the other composite let $a:P/z\to P/z$ append the identity
at $q(c)$ and retain weight $\beta$. There are natural
transformations $a\Rightarrow1$ (collapse that identity edge)
and $a\Rightarrow jr$ (identity on previous vertices and final
component $\beta$). The latter is over the identity map of the
extended type; its last square and its comma equation both read
$\beta=\beta$. These transformations prove the first specified
weak equivalence; they do not require an inverse to $\beta$.

For the second we read INDLOCAL C.7, including its two truncation
adjunctions. For paths ending at $z$ with free origin, the presheaf
$a\mapsto B(a\downarrow s)$ is the homotopy colimit of the
representables at the origins. Prepending a marked edge and its
reverse gives inverse maps up to the two explicit backtracking
transformations. Hence this presheaf inverts all marked arrows.
For any local presheaf $Q$, deleting an initial forward edge is
left adjoint to identity padding; deleting a backward edge is
right adjoint to padding. Taking opposites gives the required
initial functor in each of the two limit calculations, and in
the backward case the coefficient comparison is an equivalence
because $Q$ inverts the edge. Consequently maps from the all-type
presheaf to $Q$ are $Q(z)$. This is the universal property of
the localized representable. With every arrow marked the
localization is the fundamental infinity-groupoid of $BY$, giving
the path-space comparison stated in (C.6.1).
\end{proof}

\begin{claim}{C.7}{Proved}{Complete classification for the listed localizers}
\label{C:classification}
The following table gives necessary and sufficient conditions.
Connectivity always includes nonemptiness and one component.
\begin{center}
\begin{tabular}{lll}
\hline
Localizer & $q$ aspherical & $q$ strict, or $q$ total\\
\hline
$W_0$ & $BY$ is $1$-connected & $BY$ is $1$-connected\\
$W_1$ & $BY$ is $2$-connected & $BY$ is $2$-connected\\
$W_n$, $2\leq n<\infty$ & $BY$ is $(n+1)$-connected
  & $Y$ is a connected thin groupoid\\
$W_{\Z}$ & $BY$ is weakly contractible
  & $Y$ is a connected thin groupoid\\
$W_\infty$ & $BY$ is weakly contractible
  & $Y$ is a connected thin groupoid\\
\hline
\end{tabular}
\end{center}
\end{claim}
\begin{proof}
The path spaces in C.6 are empty off the component of $y_0$.
Within that component they have the weak type of $\Omega BY$.
Their component set is $\pi_1(BY,y_0)$ and, in their identity
component, $\pi_i\Omega BY=\pi_{i+1}BY$ for $i\geq1$.
This proves the ordinary asphericity columns for $W_n$ and
$W_\infty$. If these path spaces are integrally acyclic, they
are connected and $H_1(\Omega BY)=\pi_1(\Omega BY)=\pi_2BY$
is zero. The loop space is therefore simply connected and
acyclic, hence weakly contractible by Hurewicz in the first
possible nonzero degree. This proves the $W_{\Z}$ entry.
For $W_0,W_1$, C.5 applies both to $H_z$ and to the terminal
presheaf of $P/z$, so all three notions coincide. For the other
localizers apply C.4.
\end{proof}

\paragraph{Represented derivators and the limits of the group test.}
We checked REVLOC H.4 directly. A locally constant $C$-valued
category diagram factors through the fundamental groupoid.
If $C$ is complete, cocomplete, and not a preorder, distinct
$r,s:a\to c$ separate the product projections of $c^T$.
The regular permutation action on $c^G$ is faithful; it detects
noninjectivity of a group homomorphism. For a proper inclusion
$G<H$, projection to the coset $G$ in $c^{H/G}$ is $G$-equivariant
but not $H$-equivariant, detecting failure of fullness of
restriction. Missed or merged components are detected by,
respectively, distinct transformations and nonisomorphic objects.
This proves its $W_1$ case. In a nontrivial complete preorder all
group actions are trivial and initial/terminal values detect
exactly component bijections, giving $W_0$; for $C\simeq e$
every functor qualifies. Thus C.7 supplies the two nondegenerate
represented cases; in the last case every endpoint functor is
aspherical, strict and total, even for a disconnected $Y$.

REVLOC H5.6 is a criterion for intrinsic fibrations, a different
property from strict asphericity in XII Definition~2. Its double
coset set is the component set of the groupoid homotopy pullback,
and its displayed fiber-product groups are its isotropy groups;
the extra kernel is the image of the $\pi_2$ boundary. We checked
these identifications from the homotopy exact sequence. They
agree with the degree shift in C.7 but do not justify replacing
strict asphericity by a fibration test. H5.7 correctly distinguishes
$W_0$, the emptiness localizer, and the all-map localizer despite
identical local groups. H5.8's component-preserving qualification
is therefore essential: the relevant criterion is localization
of the full fundamental groupoid, including its components.
No extension of C.4 to an unspecified localizer is claimed.

\paragraph{Manuscript consequence and source boundary.}
At XII ms.~109, Definition~2, strict and total asphericity agree
for these full endpoint functors at every localizer listed in C.7.
For $W_\infty$ they require the categorical equivalence $Y\simeq e$;
ordinary asphericity at ms.~101--103 requires only weak
contractibility of $BY$. Thus the proposed stronger path-fibration
claim fails already for $R_0$ although its classifying space and
all endpoint slices are contractible. The length-one truncation
of PRETOT E.8 remains strictly aspherical; passage to the full
path category changes that property, exactly through the two
length-two maps in C.1.


The nearest opened classical sources are Evrard's path-category
construction and categorical versions of Theorems~A and~B
\cite{CDEvrard1975}, and Shoikhet's cocylinder and homotopy-fibrant
replacement \cite[Sections~2--3]{CDShoikhet2016}. Their endpoint
and base-change constructions concern the ordinary homotopy-fibration
property. They do not assert the strict criterion in C.4.
The exact all-type mapping-space comparison used here is the
supplied INDLOCAL C.7, audited in C.6.
